%%%%%%%%%%%%%%%%
% comm-alg.tex
%%%%%%%%%%%%%%%%


We briefly review algebraic concepts and tools used in this thesis, mainly in Chapter~\ref{chap:alg-theory}~and~\ref{chap:lowD-codes}.
There are many nice textbooks including those by Lang~\cite{Lang}, Atiyah and MacDonald~\cite{AtiyahMacDonald}, and Eisenbud~\cite{Eisenbud}.
The book by Lang is a comprehensive textbook covering a wide range of topics in abstract algebra.
The book by Atiyah and MacDonald explains commutative algebra that may look too concise, but precisely for this reason it is very useful as a reference.
Examples are rare but essential.
The book by Eisenbud is also on commutative algebra and is extensive.
It covers more material than Atiyah-MacDonald.
In particular, our summary of Gr\"obner basis follows Eisenbud.
The chapter on Gr\"obner basis appears in the middle of the book,
but is relatively self-contained and elementary.
Here, we will omit many proofs and not try to be fully rigorous.
We explain theorems to the point where intuition can be developed.
Rigorous proofs can be found in one of the three books.



We start by recalling definitions for abelian groups.
An abelian group $G$ with the identity element denoted by $0$
is a set with an operation $+ : G \times G \to G$ such that
$g + g' = g' + g$ and $0 + g = g$.
It is required for $G$ to have inverses of $g$ denoted by $-g$ such that $g + (-g) = 0$.
$n$-fold sum of $g$ is simply denoted as $ng$, where $n \in \ZZ$.
Given two abelian groups $G$ and $H$, we can form a direct sum $G \oplus H$.
It is the set of all tuples $(g,h)$, where $ g \in G,~ h \in H$,
and the group operation $+$ is defined as $(g,h) + (g',h') = (g+g',h+h')$.
We can form a {\bf direct sum} $\bigoplus_\alpha G_\alpha$ of arbitrary family $\{ G_\alpha \}$ of groups.
It is the set of all indexed collections of group elements $(g_\alpha)$ 
where only {\em finitely many} $g_\alpha$ are nonzero.
The group operation is again defined component-wise.
Thus, any element in the direct sum is a sum of finitely many $g_\alpha \in G_\alpha$.
A sum of two abelian groups can be defined if they are subgroups of a parent group.
If $A,B \le C$ are subgroups, the {\bf sum} $A + B$ is the group of all elements of $C$ of form $a+b$ where $a \in A$ and $b \in B$.
Note that $A+B \cong A \oplus B$ if and only if $A \cap B = 0$.
Given a subgroup $N \le G$, we can form a quotient group $G/N$,
the set of all equivalent classes under the equivalence relation $[g]=[g']$ iff $g-g' \in N$.
In commutative algebra, almost everything is an abelian group.
On top of the abelian (additive) group structure, a new ``multiplication'' is added.


\section{Rings and homomorphisms}

The set of integers $\ldots, -2,-1,0,1,\ldots$ admits two operations, addition and multiplication.
There is 0 that has no effect under addition, and 1 that has no effect under multiplication.
One can always undo the addition because one can subtract a number.
However, the multiplication is not invertible within the set of integers
because fractions are not integers.
One convenient thing is that the multiplication does not care about the order.
A commutative ring is an abstraction of this structure.
It is a set, in which one can add and subtract. A multiplication exists but is not in general invertible.
An additive identity 0 exists, and a multiplicative identity 1 exists.
The distribution law $a(b+c) = ab + ac$ is assumed,
and the multiplication is commutative $ab = ba$.
A ring $R$ can consists of a single element, in which case $R$ is called a zero ring, if and only if $0=1$.
Indeed, if $a \in R$ and $1 = 0 \in R$, then $a = a \cdot 1 = a \cdot 0 = a \cdot ( 0 + 0 ) = a \cdot (1+1) = a + a = 0$.
Examples of rings are abundant: The set of all integers, the set of all complex numbers, the set of all square diagonal matrices of  a fixed size,
the set of polynomials, the set of all differentiable functions on a real line, the set of all continuous real-valued functions on a manifold, etc.
Is the set of all even integers a ring? 
No. Some authors define rings to include this case where the multiplicative identity $1$ is not provided,
but we avoid this case. Any ring is with 1.
Note that $1$ is unique; if $1'$ is also a multiplicative identity, then $1 = 1 \cdot 1' = 1'$.
The same is true for $0$.

A ring is always understood in terms of relations with other rings.
Given two rings $A$ and $B$ we consider a restricted class of maps between them.
That is, we require that the map obeys the ring structure of the rings.
$f : A \to B$ is a {\bf homomorphism} if $f(a+b) = f(a) + f(b)$ and $f(ab) = f(a)f(b)$ for any $a,b \in R$.
In addition, we assume $f(0) = 0$ and $f(1) = 1$.
(``morph'' means ``shape.'') 
The $+$ or the omitted $\cdot$ between $a$ and $b$ in $ab$ on the left-hand side
are the operations defined in $A$, whereas those in the right-hand side are in $B$.
The {\bf image} of a homomorphism $f$ is the subset of $B$ written as $f(A)$ defined by $\{ f(a) ~|~ a \in A \}$.
Is the image of a homomorphism a ring? Yes.

There is no point to speak of a map between two rings $A,B$ that is not a homomorphism.
If we are going to ignore the ring structure, we would rather say the map between the ``sets'' $A,B$.
We will simply say a {\bf map between rings} to mean a homomorphism.
We note more terminologies: An {\bf endomorphism} is a map from a ring into itself.
An {\bf isomorphism} is a map between two rings with a unique inverse.
An {\bf automorphism} is an isomorphism that is an endomorphism.

The ring of integers is so primitive in the following sense.
Let $A$ be an arbitrary ring. Consider a map $f : \ZZ \to A$.
$f(n) = \sum_{i=1}^{n} f(1)$ and $f(-n) = \sum_{i=1}^n (-f(1))$ where $n > 0$.
But, $f(1)=1$ is the unique multiplicative identity.
Therefore, $f$ is completely determined, though we just required $f$ be a homomorphism;
there is a unique nonzero map from $\ZZ$ into any ring.
How many endomorphisms are there for $\ZZ$?

The {\bf kernel} of a map (homomorphism!) $f$ is the subset of $A$ written as $\ker f$
defined by $\{ a \in A ~|~ f(a) = 0 \}$.
It is easy to see that the kernel is closed under the addition and multiplication.
Here, the closeness means that the result of the operation using two elements in a subset lies in the subset.
(It is pointless to speak of the closedness of an operation without reference to a subset.)
There is one more important property as we discuss below.


\section{Ideals and modules}

The kernel $I$ of a map $f$ between rings $R \to S$ has the following property:
\begin{equation}
 \forall r \in R,~ \forall x \in I~:~~ r x \in I
\label{cond-ideal}
\end{equation}
This is easily verified since $f(r x) = f(r)f(x) = f(r) \cdot 0 = 0$.
When a subset $I$ of a ring $R$ is closed under the multiplication and addition, and satisfies the above property,
we call $I$ to be an {\bf ideal} of $R$.
For example, in $\ZZ$, the set of all even numbers is an ideal denoted as $(2)$.
The property \eqref{cond-ideal} is trivial, because it reads a multiple of an even number is even.
This ideal is the kernel of the map $\ZZ \to \ZZ/(2)$,
where the latter is the ring of integers modulo 2.
Remember that there is a unique nonzero map from $\ZZ$ to any ring.
In fact, any ideal arises in this way: An ideal is the kernel of a ring homomorphism.

To understand this, we need to formalize how to construct {\bf quotient rings} or {\bf factor rings}.
Let $R$ be a ring and $I$ be an ideal.
They are both abelian groups; they are closed under the addition with identity $0$,
and contains additive inverses, the minus elements.
The quotient ring $R/I$ as a set is the same as the quotient group as an abelian group.
$R/I$ is the family of equivalence classes under the equivalence relation that
$[a] = [b] \in R/I$ iff $a-b \in I$.
The multiplication in $R/I$ is as expected: $[a] \cdot [b] = [ab]$.
Using the property $(*)$ one can verify that it is well-defined.
Consider a map $R \to R/I$ defined by $a \mapsto [a]$.
It should be straightforward that it is a ring homomorphism.
What is the kernel? Precisely, $I$.
Now it is an almost tautology that an ideal is a kernel of a ring homomorphism.
Often, we just write $a$ in place of the equivalence class notation $[a]$.
In this lazy notation, ``the map $R \to R/I$ is defined by $a \mapsto a$.''
This map is so important that it has its own name: {\bf canonical map} or {\bf quotient map}.


In leisurely words, an ideal extends what we treat as zeros;
any element of $I$ is zero in $R/I$.
Zero plus zero or zero times zero must be zero, so $I$ is closed under the addition and multiplication.
Zero times any element must be zero, so the property \eqref{cond-ideal} should hold.
How do we specify an ideal in a real calculation?
Consider a polynomial ring $\QQ[x]$, the set of all polynomials (of finitely many terms) in $x$ with coefficients in the rational numbers.
Suppose we decided that $x$ must be equal to $\half$.
In other words, we decided that $x - \half$ must be ``zero'' in $\QQ[x]/I$.
Then, as zero times anything must be zero, any element $(x-\half)\cdot f(x)$ must be zero in $\QQ[x]/I$, too;
it must be an element of the ideal we are defining.
We require no more elements are identified as zeros.
That is, we define $I = \{ f(x)(x-\half) ~|~ f(x) \in \QQ[x] \}$ as our ideal,
and carry out any computation in the quotient ring $\QQ[x]/I$.
We write $I$ as $(x-\half)$ by putting the {\bf generator} inside the parenthesis.
Computing in $\QQ[x]/I$ is the same as computing in $\QQ[x]$ and evaluate the polynomial at $x=\half$.


In general, if we write and ideal $I$ of $R$ as $(a,b,c)$,
then we mean
\[
 I = \{ ax + by + cz ~|~ x,y,z \in R \} .
\]
$I$ is said to be {\bf generated by $a,b,c$}.
A quick exercise: What is $\QQ[x]/(x-\half, x-1)$~? It is a zero ring.
Since $x-\half$ and $x-1$ are ``zeros,'' their difference $\half$ is zero.
Zero times 2 is zero, so $1$ is zero. Therefore, everything is zero.
We proved an ideal equality $(x-\half, x-1) = (1)$.
The whole ring $R$ viewed as an ideal is called the {\bf unit ideal} denoted by $(1)$.
In fact, when there is any invertible element in an ideal, it is the unit ideal.
For this reason, an invertible element in a ring is called a {\bf unit}.
Note that computing the minimal set of generators is in general a hard problem.
For instance, I do not know any algorithmic answer to questions like ``can this ideal be generated by three elements?''
However, the Gr\"obner basis gives a \emph{canonical} set of generators for an ideal of a polynomial ring,
and we can algorithmically compare two ideals.

\subsection*{Prime, maximal}

There is a very important class of ideals that generalizes prime numbers in $\ZZ$.
A prime number $p$ has a property that if $p$ divides a product of two integers $ab$
then $p$ divides either $a$ or $b$. A {\bf prime}%
\footnote{Do not use ``primary'' in place of ``prime.'' The adjective ``primary'' has a slightly different technical meaning.}
ideal $\pp$ of $R$ is an ideal not equal to $(1)$ such that
\begin{equation}
ab \in \pp \text{ implies } a \in \pp \text{ or } b \in \pp \text{ for any } a,b \in R.
\label{cond-prime}
\end{equation}
The condition is rephrased as $a \notin \pp$ and $b \notin \pp$ imply $ab \notin \pp$.
Is the ideal $(0)=\{0\}$ prime? It depends. In $\ZZ$ the zero ideal is prime because the product of two nonzero integers is nonzero.
In the polynomial ring $\QQ[x]$, $(0)$ is prime because the degree of nonzero polynomial does not decrease under any nonzero multiplication. 
If any nonzero elements of a nonzero ring $R$ has a multiplicative inverse, in which case $R$ is called a {\bf field},
then $(0)$ is prime because $ab = 0$ means $a = 0$ or $b = 0$.
However, in the ring of diagonal $2 \times 2$ matrices, 
$(0)$ is not prime because nonzero matrices $\begin{pmatrix} 1 & 0 \\ 0 & 0 \end{pmatrix}$
and $\begin{pmatrix} 0 & 0 \\ 0 & 1 \end{pmatrix}$ multiply to zero.
The ring in which $(0)$ is prime has a special name, {\bf (integral) domain}.
It is easy to verify that $R/\pp$ is an integral domain if and only if $\pp$ is prime,
applying the picture that $\pp$ defines zeros in $R/\pp$.


There is one more important property of prime ideals.
Let $f : A \to B$ be a map between two rings.
{\em If $\pp \subseteq B$ is a prime ideal, then $I = f^{-1}(\pp) \subseteq A$ is prime.}
Proof: $aa' \in I \Rightarrow f(a)f(a') \in \pp \Rightarrow f(a) \in \pp \vee f(a') \in \pp \Rightarrow a \in I \vee a' \in I$.

A subclass of prime ideals consists of maximal ideals.
A {\bf maximal ideal $\mm \neq (1)$} is defined by the maximal property:
\begin{equation}
 \mm \subsetneq \mm' \text{ implies } \mm' = (1) \text{ for any ideal } \mm'.
\end{equation}
It is a priori not vivid why maximal ideals are prime.
But it is easy to see. Consider $R/\mm$. If $a \in R/\mm$ is nonzero, that is $a \notin \mm$, then $I = (\mm, a) \supsetneq \mm$ and therefore $I = (1)$, which means there is an element $b \in R$ such that $b a + m = 1$ for some $m \in \mm$. By the canonical map, $b$ maps to a multiplicative inverse of $a$ in $R/\mm$. 
(The converse is also true. {\em $\mm \neq (1)$ is a maximal ideal if and only if $R/\mm$ is a field}.)
In other words, $R/\mm$ is a field, and therefore an integral domain. In particular, $\mm$ is prime.

\subsection*{Modules}

Ideals admit another viewpoint.
Let us forget the multiplication within an ideal $I$,
and treat $I$ as a separate set from the mother ring $R$.
It is still an abelian group, and satisfies \eqref{cond-ideal}.
The condition $(*)$ looks very similar to the scalar multiplication on vector spaces.
Indeed, consider a direct sum $R \oplus R$ of abelian groups,
the set of all tuples $(r,r')$ where $r,r' \in R$ are any elements.
There is an operation on $R \oplus R$ similar to \eqref{cond-ideal};
we can define $r \cdot (s,s')$ to be $(rs, rs')$,
similar to the scalar multiplication for a vector space.
Indeed, a two-dimensional vector space is precisely obtained in this way by setting $R = \mathbb{Q}, \mathbb{R}, \mathbb{C}$, etc.


Let us define an abstract notion.
Let $M$ be an abelian group with a {\em bilinear} operation $\cdot$ such that
\begin{equation}
 \forall r \in R,~ \forall m \in M~:~~ r \cdot m \in M .
\label{cond-module}
\end{equation}
The ``multiplication'' $\cdot$ here is not the same thing as the multiplication within the ring.
We assume expected formulas $(rs) \cdot m = r \cdot (s \cdot m)$, and simply write $(rs) \cdot m = rsm$ where $r,s \in R$ and $m \in M$.
We call $M$ an {\bf $R$-module} or a {\bf module over $R$}.
It is general than the notion of vector spaces.
A vector space is a module over a {\bf field}, a commutative ring in which multiplication by a nonzero element has an inverse.
An ideal is a subset of $R$ such that it is a module over $R$.%
%%%
\footnote{
Although an ideal is a valid module, often it is treated differently than a module.
One should sometimes be careful to apply things defined for modules,
especially when one reads the dimension theory of Eisenbud~\cite{Eisenbud}.}
%%%
Note that $R$ itself is an $R$-module via the multiplication within $R$.
The above example $R \oplus R$ is an $R$-module.


For an $R$-module $M$ if there exist finitely many elements $m_1,\ldots,m_n \in M$ such that
\begin{equation}
 M = \{ r_1 m_1 + \cdots + r_n m_n ~|~ r_1,\ldots, r_n \in R \},
\end{equation}
then we say $M$ is {\bf finitely generated}.
All of our modules in the thesis are finitely generated.
An ideal is finitely generated if it is finitely generated as a module.

There is a confusingly similar terminology that one must distinguish.
Suppose $A$ is a subring of $B$. That is, $A$ is a ring by itself and contained in a bigger ring $B$.
Or slightly more generally, suppose we are given a ring map $A \to B$.
The subring case is precisely when the map is an inclusion.
We say $B$ is a {\bf finitely generated $A$-algebra}
if there exists finitely many elements $b_1, \ldots, b_n \in B$ such that
any element of $b$ can be written as a polynomial in $b_1,\ldots, b_n$ with coefficients in the image of $A$.
In this case, $B$ is sometimes written as $B = A[b_1,\ldots, b_n]$.
A typical situation is when $A$ is a field such as $\QQ, \mathbb{C}$ and $B$ is a polynomial ring over $A$.
For example, $B = \QQ[x,y]$.
The reason it is a confusing terminology is because $B$ is not necessarily finitely generated $A$-module.
$\QQ[x]$, a finitely generated $\QQ$-algebra, has infinitely many generators $\{ 1, x, x^2, \ldots \}$ as a $\QQ$-module.
Is $\QQ$ a finitely generated $\ZZ$-algebra? 
No, because multiplying by integers cannot produce large denominators.


As the rings are understood in relation to others,
the modules should be understood via maps.
We required the ring homomorphisms to preserve the defining operations of the rings.
The same is true for the module maps. Given two modules $M$ and $N$ over $R$,
we define an {\bf $R$-linear map} or {\bf $R$-module homomorphism} $f : M \to N$ to satisfy
\begin{equation}
f(m+m') = f(m)+f(m')\quad \text{ and } \quad f(r \cdot m) = r \cdot f(m) \text{ for any } r \in R,~ m, m' \in M.
\end{equation}
Note that the $\cdot$ on the left-hand side is the operation \eqref{cond-module} for $M$
whereas that on the right-hand side is the operation for $N$.%
\footnote{
In group representation theory,
the $R$-linear maps are called equivariant maps,
where $R$ is the group algebra which may not be commutative.
The representation space is a module, 
the subrepresentation space is a submodule, and
the irreducible representation space is a simple submodule.}
The notions of kernel, image, endomorphism, automorphism, and isomorphism apply to module maps, too.
A simple exercise: Let $A,B$ be two $R$-algebras; there are ring maps $R \to A$ and $R \to B$.
The two algebras are naturally $R$-modules, when a map $A \to B$ becomes $R$-linear?
Answer: It becomes $R$-linear when the following diagram commutes.
\[
 \xymatrix@!0{
A \ar[rr]  & & B \\
           & R\ar[ul]\ar[ur] &
}
\]


\subsection*{Free}

An $R$-module isomorphic to $\bigoplus_\alpha R$ is called a {\bf free module}.
When there are finitely many summands, it is called {\bf finitely generated free module}.
It is the most convenient type of modules.
In particular, a module map between finitely generated free $R$-modules
is simply given by a matrix with entries in $R$, just as linear map between finite dimensional vector space
can be described by a matrix with entries in a field.
Note that elements of a finitely generated free $R$-module $M = R^{\oplus n}$ can be expressed by column matrices.
Let $e_i$ ($i=1,\ldots,n$) denote the canonical basis column matrices.
(We could say ``column vectors'' instead of column matrices. However, an element of a module is not a vector in general.)
A map $f$ from $M$ to any module is specified if we specify the image $f(e_i)$
because the image of other elements $a_1 e_1 + \cdots a_n e_n$ is must be $a_1 f(e_1) + \cdots + a_n f(e_n)$
by $R$-linearity.
Writing $f(e_i)$ in columns, we have a matrix representation of $f$.
An un-redundant set of generators of a free module is called a {\bf basis}.
The cardinality of a basis is called {\bf rank}.
(One can show that rank is independent of the choice of a basis.)
Only for free modules can we speak of bases.
The crucial difference between general modules and vector spaces is that
a module is in general not free, while a vector space, a module over a field, is always free.
That nonzero elements are invertible makes such a huge difference.


Note that any module can be described by free modules.
Take a {\bf generating} set $\{ m_\alpha \}$ of a module $M$; 
any element of $M$ is a {\em finite} $R$-linear combination $\sum_i r_{\alpha_i} m_{\alpha_i}$.
A trivial and useless choice would be to take whole $M$ as a generating set.
Let $F$ be a free $R$-module whose rank is the same as the cardinality of the generating set,
i.e., there is a surjective module map $F \to M$.
The kernel $N$ is a submodule of $F$, not necessarily free, and $M \cong F / N$.
One can carry the same process for the module $N$.
That is, one finds a free module $F'$ such that $\phi : F' \to N$ is a surjection.
Now $M$ is expressed as $M \cong F / \im \phi$.
This is conceptually important observation, but not too useful
because we do not have any control over the ranks of $F$ and $F'$.
We need some finiteness conditions.
An $R$-module $M$ is said to be {\bf finitely presented} 
if there is a map $\phi : F' \to F$ between finitely generated free modules
such that $M \cong F / \im \phi$. The latter expression $F/\im \phi$ is often abbreviated as $\coker \phi$.
The map $\phi$ is called a {\bf finite presentation} of $M$.
As we have seen above, $\phi$ is a matrix with entries in $R$,	
and $M$ is expressed by a single matrix $\phi$.
In case of a finite dimensional vector space, $\phi$ can always be brought to a diagonal matrix
with entries of $0$ or $1$, after basis change of $F$ and $F'$.
We will discuss more on a finite presentation in Section~\ref{sec:det-ideal}.


\subsection*{Noetherian}

A technically very important and convenient adjective is Noetherian.
It is as important as vector spaces having finite dimensions.
A {\bf module $M$ is Noetherian} if every increasing sequence of submodules is stationary,
i.e., if $M_1 \le M_2 \le \cdots \le M_n \le \cdots$ is a sequence of submodules of $M$,
then for all sufficiently large $n$ one has $M_n = M_{n+1}$.
Often this condition is referred to as the {\bf ascending chain condition} or {\bf a.c.c.}
{\em A module $M$ is Noetherian if and only if any submodule is finitely generated.}
If a submodule cannot be generated by finitely many elements, 
one can construct a strictly increasing infinite sequence of submodules
using an infinite subset of generators.
Conversely, if any submodule of $M$ is finitely generated,
then the union $\cup_{i=1}^\infty M_i$ of any increasing sequence of submodules $M_i$ of $M$
is also finitely generated, say, by $m_1,\ldots,m_r$.
Since $m_j$ is contained in some $M_{j'}$, there must be some $k$ such that $m_1,\ldots,m_r \in M_k$.
It follows that $M_k = M_{k'}$ for all $k' \ge k$.
A Noetherian module over a field is just a finite dimensional vector space.
A {\bf Noetherian ring} $R$ is a ring that is Noetherian as an $R$-module,
i.e., the a.c.c is satisfied with respect to the ideals of $R$.
In a Noetherian ring, any ideal is finitely generated.


Being Noetherian is preserved in many cases.
Loosely speaking, it says that a module carries a finite amount of data.
A module with a finite amount data manipulated finitely many times would still have a finite amount of data.
The following are facts:
\begin{itemize}
\item {\em A homomorphic image of Noetherian ring (module) is a Noetherian ring (module).}
\item {\em A submodule of Noetherian module is Noetherian.}
\item {\em A finitely generated algebra over a Noetherian ring is a Noetherian ring.}
\item In particular, {\em a polynomial ring over a field with finitely many variables is a Noetherian ring.}
\item {\em A finitely generated module over a Noetherian ring is Noetherian.}
\end{itemize}
There are more to mention about being Noetherian using tensor products and localization.
See Section~\ref{sec:localization}.

Remark that {\em a finitely generated module $M$ over a Noetherian ring $R$ always admits a finite presentation $\phi : F' \to F$.}
Since $M$ is finitely generated, the module $F$ is of finite rank.
$F$ is Noetherian, and therefore, $\im \phi$ is finitely generated, which means $F'$ can be taken to be of finite rank.
The ``matrix'' $\phi$ is of finite size.

Is the ring of all differentiable functions $\mathbb{R} \to \mathbb{R}$ Noetherian? No.
Is the ring of all trigonometric functions $\mathbb{R} \to \mathbb{R}$ of period $1/n$, where $n$ are positive integers, Noetherian? Yes.

\section{Gr\"obner basis}

Gr\"obner basis is a special generating set for ideals and modules.
It provides canonical presentations of modules and ideals, from which any concrete computational commutative algebra is built.
In all notions in the previous section,
I cannot imagine any systematic way to compute things concretely without Gr\"obner basis.
It is theoretically important too because it tells what is actually constructible.
A nice application of the Gr\"obner basis is a sharp version of Hilbert syzygy theorem~\cite[Corollary~15.11]{Eisenbud}.
In this section we {\em assume the (base) ring is a polynomial ring over a field $\FF$ with finitely many variables};
$R = \FF[x_1,\ldots,x_n]$.
Since $R$ is Noetherian, for any ideal there is a finite set of generators.

How do ideals look like in $\FF$? There are only $(0)$ and $(1)$ because any nonzero element is a unit.
How about $\FF[x]$? It is only slightly more complicated.
By the Euclidean algorithm, the $\gcd$ of two polynomials $f(x),~g(x)$ can be expressed as
\[
 \gcd( f(x), g(x) ) = a(x) f(x) + b(x) g(x) \in (f(x),g(x))
\]
An ideal generated by $f(x)$ and $g(x)$ is thus the same as an ideal generated by a single element $\gcd(f(x),g(x))$, 
in which case the ideal is called {\bf principal}.
By induction, one can always reduce the number of generators of an ideal in $\FF[x]$ if it is greater than $1$.
Since $\FF[x]$ is Noetherian, this is enough to imply that {\em any ideal of $\FF[x]$ is generated by a single element,}
i.e., {\em any ideal of $\FF[x]$ is principal}.
The principal generator $p(x)$ is important, not only because it is simple, but also because it gives a criterion
whether an arbitrary polynomial $h(x)$ is contained in $(p(x))$:
Divide the $h(x)$ by $p(x)$. The remainder is zero if and only if $h(x) \in (p(x))$.
Put differently, we can find a canonical representative in the quotient ring $\FF[x]/(p(x))$
to be the remainder $r(x)$  where $\deg r < \deg p$.
The $p(x)$ is the Gr\"obner basis for the ideal $(p(x))$ in $\FF[x]$.

In case of two or more variables, it is no more true that any ideal is principal.
Let us examine the {\bf division algorithm}.
Let $f(x)$ be a dividend and $p(x)$ be a divisor.
We first compare degrees of them.
If $\deg f < \deg p$, then the division is completed, and $f$ is the remainder.
Otherwise, we match the leading coefficients and then subtract a multiple $a(x)p(x)$ from $f(x)$.
$a(x)$ is a monomial such as $3x^2$.
The purpose is of course to have $\deg (f-ap) < \deg f$.
And then we iterate.
Essential is a total ordering among terms such that the ordering is preserved under multiplication by a monomial.
Now we consider a problem of computing the remainder of $f$ modulo $G = \{ g_1, \ldots, g_k \}$.
Equipped with the ordering, given a dividend $f$ and a set $G$ of divisors,
(Step-1) one should be able to match the leading term and kill it,
thereby ``reduce'' the leading term of the dividend.
(Step-2) One stops when the reduction becomes impossible.
A question arises.
During Step-1, there would be many possible choices among the divisors from $G$.
How do we guarantee that the remainder is independent of the choices?
This is a sound question, and the answer is no in general.
Even in the one variable case, the answer is no in general.
Consider $f = x^2 -1,~ g = x^3 -1$ in $\QQ[x]$.
The division by the set $\{f,g\}$ of a dividend $x^3$ gives $1$
if we kill the leading term by $x^3-1$,
or $x$ by $x^2-1$.

We formulate the problem as finding a generating set of an ideal $I$ 
that gives rise to a unique remainder under the division.
A solution is that we include more and more elements of $I$ into a generating set $G$
so that any leading term $l$ of a dividend can be removed 
where $l$ appears as a leading term of some elements of $I$.
In the one variable case, it is enough to include the $\gcd$ of all (finitely many) generators of $I$.
We have given enough motivation and technical problems to smooth out.

\subsection*{Definitions}

A monomial is a product of variables.
A {\bf monomial order} on the set of all monomials of $R = \FF[x_1,\ldots,x_n]$
is a total ordering $\succ$, under which any two monomials are comparable, such that
\begin{equation}
 x_i m \succ x_i m' \succ m' \quad \text{whenever} \quad m \succ m'
 \label{cond-TermOrder}
\end{equation}
for any monomial $m,m'$ and any variable $x_i$.
Since a monomial $x_1^{a_1} \cdots x_n^{a_n}$ is uniquely given by a $n$-tuple of nonnegative integers $(a_i) = (a_1,\ldots,a_n)$,
a monomial order is a total order on the hyper-octant of $\ZZ^n$ of nonnegative coordinates.
Note that under any monomial order, $1$ is the least monomial.

Two examples at least are important. The first one is the {\bf lexicographic order},
under which $(a_i) \succ (b_i)$ if and only if $a_i > b_i$ for the least $i$ such that $a_i \neq b_i$.
For example,
\[
 x_1 \succ_\text{lex} x_2^5 \succ_\text{lex} x_2 x_3^{100} .
\]
The second example is the {\bf degree reverse lexicographic order},
under which $(a_i) \succ (b_i)$ if and only if `$\sum a_i > \sum b_i$' 
or `$\sum a_i = \sum b_i$ and $a_i < b_i$ for the largest $i$ such that $a_i \neq b_i$.'
For example,
\begin{align*}
 x_1 \succ_\text{revlex} x_2 \succ_\text{revlex} \cdots \succ_\text{revlex} x_n \\
 x_3^{100} \succ_\text{revlex} x_2^{50} \succ_\text{revlex} x_1^{49} x_3 .
\end{align*}
The degree reverse lexicographic order is the most commonly used in computer software.


{\em Any monomial order is a well-order, i.e., every set of monomials has a minimal element.}
This statement is at the core of any finiteness proof regarding Gr\"obner basis.
The proof is very simple. Let $X$ be any set of monomials of $R$.
The submodule (ideal) generated by $X$ is finitely generated by $G$ because $R$ is Noetherian.
It means $X$ consists of multiples of a finitely many monomials of $G$.
Therefore, the minimal element of $G$, a finite set, is the minimal element of $X$.
The statement can be rephrased as {\em any decreasing sequence of monomials is stationary}
or as {\em any strictly decreasing sequence of monomials is finite.}


With a monomial order, we can define the {\bf leading term} or {\bf initial term} of a polynomial $f$
which is the greatest term of $f$ with respect to the monomial order.
(We are distinguishing ``term'' and ``monomial.'' A monomial is a product of variables only, and a term is $\FF$-multiple of a monomial.)
Let us denote the leading term of $f$ by $\lt(f)$.
Given any set $S$ of $R$, let the ideal generated by all leading terms of element of $S$ be denoted by $\lt(S)$.
Now, a {\bf Gr\"obner basis} $G$ of $I$ is a generating set of $I$ such that $\lt(I) = \lt(G)$.
Since $R$ is Noetherian, $\lt(I)$ is finitely generated.
Therefore, $G$ can be chosen to be a finite set.

\subsection*{Buchberger criterion}

{\em The division by a Gr\"obner basis results in a unique remainder.}
In particular, a polynomial is in $I$ if and only if the division by a Gr\"obner basis yields the zero remainder.
To see this, let $f$ be an arbitrary polynomial.
Let $r$ be a remainder obtained by killing large terms of $f$ by elements of $G$.
Let $r'$ be another remainder. We must show $r = r'$.
It is clear that $r-r'$ belongs to $I$. If nonzero, the leading term of $r-r'$ belongs to $\lt(I) = \lt(G)$.
However, $r$ or $r'$ has by construction no leading term that belongs to $\lt(G)$.
Therefore, $r-r'$ must be zero.


It is not clear yet how $G$ can be computed from given generators of $I$,
though we know a solution in the case where the generators of $I$ are in a single variable.
A strategy is hinted from one-variable case:
\begin{itemize}
\item[(1)] Start with any generating set $S$ of $I$. 
\item[(2)] Try to produce new polynomials whose initial terms are not contained in $\lt(S)$.
\item[(3)] Update $S$ by adjoining the new polynomials.
\item[(4)] Iterate.
\end{itemize}
It will end after finitely many iterations because the new initial terms form a descending sequence of monomials.
It remains to find an effective method to produce new initial terms.
Let us work with an example first.
Consider $I = (x^2,xy+y^2) \subset \QQ[x,y]$.
We will compute Gr\"obner basis with respect to two monomial orders.
The first one is the lexicographic order under which $x \succ y$.
We have $\lt(x^2) = x^2$ and $\lt(xy+y^2) = xy$.
Neither of them is a multiple of the other.
However, $y(x^2) - x(xy+y^2) = -xy^2$ has the leading term divisible by $xy$.
This combination is called ``S-polynomial'' of $x^2$ and $xy+y^2$.
So, $y(x^2) - x(xy+y^2) + y(xy+y^2) = y^3$.
We see that $y^3$ is a new initial term of $I$.
The generating set is expanded as $G = \{ x^2, xy + y^2, y^3 \}$.
We then again form S-polynomials using any two among $G$,
but fail to produce any new initial term. Indeed, $G$ is a Gr\"obner basis of $I$.
The second monomial order is the lexicographic order under which $x \prec y$.
Then, $\lt(x^2) = x^2$ and $\lt(xy+y^2) = y^2$.
We form an S-polynomial, $y^2(x^2)-x^2(xy+y^2) = -x^3 y$.
It is a multiple of $x^2$, and we do not get anything new.
Indeed, $\{x^2, xy+y^2\}$ is a Gr\"obner basis.

Given two polynomials $f,g$ let us fix the leading coefficient of $\gcd$ to be $1$.
The {\bf S-polynomial} of two polynomials $f$ and $g$ is the polynomial
\[
 \sigma(f,g) = \frac{\lt(g)}{\gcd(\lt(f),\lt(g))} f - \frac{\lt(f)}{\gcd(\lt(f),\lt(g))} g .
\]
Dividing an S-polynomial by a set of generators will give a non-unique remainder,
but we can hope that the calculation would find a new initial term.
Buchberger showed that {\em the process is sufficient to find all initial terms of an ideal}.
The {\bf Buchberger algorithm} to find a Gr\"obner basis is given by the above ``algorithm,''
where the step (2) is now well-defined by (2$'$):
\begin{itemize}
 \item [(2$'$)] For each pair $f,g$ of polynomials of $S$, compute $\sigma(f,g)$ and its remainder $r$ after division by $S$.
\[
\frac{\lt(g)}{\gcd(\lt(f),\lt(g))} f - \frac{\lt(f)}{\gcd(\lt(f),\lt(g))} g = \sum_k h_k s_k  + r
\]
where $s_k \in S$, $h_k \in R$, and $\lt(r) \notin \lt(S)$ or $r=0$.
\end{itemize}
In other words, the algorithm terminates with a correct answer, a Gr\"obner basis $G$,
if the remainder of any S-polynomial computed from a pair of polynomials in $G$ is zero.
This is called the {\bf Buchberger criterion}.
It follows that if $G$ is a Gr\"obner basis for the ideal $(G)$,
then any subset $G' \subseteq G$ is a Gr\"obner basis for the ideal $(G')$.



A Gr\"obner basis $G={f_1,\ldots,f_t}$ is {\bf reduced}
if any term of $f_i$ is not divisible by any $\lt(f_j)$ where $i \neq j$,
and the leading coefficients of all $f_i$ are $1$.
The reduced Gr\"obner basis can be obtained from any Gr\"obner basis.
{\em Given a monomial order, a reduced Gr\"obner basis is unique for an ideal.}
Therefore, {\em two ideals are the same if and only if their reduced Gr\"obner bases are the same.}


\subsection*{Applications}

An immediate application is to decide whether 
a set of polynomial equations $\{f_i(x_1,\ldots,x_n) = 0 \}$ admits a common solution.
If one obtains $1 = 0$ while manipulating $f_i$, then certainly there is no solution,
and this is the only case. That is, the ideal $I$ generated by $f_i$ contains $1$,
i.e., $I$ is the unit ideal,
if and only if there is no common solution to $f_i = 0$.
Since the ideal membership can be tested by the division algorithm using a Gr\"obner basis,
we can algorithmically answer this question by computing a Gr\"obner basis
in any monomial order and looking for $1$.

Note that known algorithms for Gr\"obner basis is not efficient in a computation complexity sense.
Any significant improvement seems impossible because if one can compute Gr\"obner basis efficiently,
then one can also solve, for example, Boolean satisfiability problem (SAT) efficiently.
A SAT is a decision problem that asks whether there exists an assignment to Boolean variables $x_1,\ldots,x_n$
such that a given formula using AND, OR, and parenthesis evaluates to $1$.
Any SAT can be formulated as a finding a solution to a system of polynomial equations.
Being Boolean can be expressed as $x_i(x_i-1) = 0$.
The OR of two binary variables $x,y$ can be expressed as $1-(1-x)(1-y)$.
The AND of two binary variables $x,y$ is the product $xy$.
In particular, the 3-SAT, the SAT where all the expressions are given by conjunctive normal form with 3 variables per clause,
is equivalent to finding a solution to systems of polynomial equations, each of which has degree at most 3.
The 3-SAT is known to be NP-complete.
However, finding Gr\"obner basis for a fixed number of variables is efficient,
i.e., {\em the number of calculation steps increases as a polynomial in the degree of the generating polynomials}.


We can also compute the intersection of an ideal with a subring.
The foundational case is when the subring is given by $T = \FF[x_1,\ldots,x_n]$
contained in $R = \FF[x_1,\ldots,x_n,y_1,\ldots,y_m]$.
If an ideal $I$ of $R$ is given, we wish to compute the intersection $T \cap I$,
which is an ideal of $T$.
We should use an {\bf elimination monomial order} under which $\lt(f) \in T$ implies $f \in T$.
The lexicographic order in which $x_a \prec y_b$ is an elimination order.
The following statement is true.
{\em If a Gr\"obner basis $G=\{f_1,\ldots,f_u,g_1,\ldots,g_v\}$ of $I$ under elimination order
is such that $f_i$ do not involve any variable $y_b~(b=1,\ldots,m)$ but $g_j$ do,
then $G' = \{f_1,\ldots,f_u\}$ is a Gr\"obner basis of $I \cap T$.}
The ideal $J$ generated by $G'$ is certainly contained in $I \cap T$.
It is clear that $G'$ is a Gr\"obner basis for $J$ because it satisfies the Buchberger criterion.
If $J \subsetneq I \cap T$, then there would be $f \in (I \cap T) \setminus J$ with the least leading term.
In particular, $\lt(f) \notin \lt(J)$.
Since $G$ is a Gr\"obner basis, $\lt(f) \in \lt(G)$.
Since $\lt(g_j)$ involve $y_b$ due to the elimination order, but $\lt(f) \in T$,
we must have $\lt(f) \in \lt(G') = \lt(J)$, a contradiction.


More generally, suppose we are given with a subring $T \subseteq \FF[y_1,\ldots,y_m]$
generated by some polynomials $h_i (y_1,\ldots,y_m)$ ($i=1,\ldots,n$) over $\FF$.
We wish to find $T \cap I$ for an ideal $I$ of $\FF[y_1,\ldots,y_m]$.
Introduce new variables $x_1,\ldots,x_n$, and consider 
\[
T' = \FF[x_1, \ldots, x_n] \xrightarrow{\phi} \FF[y_1,\ldots,y_m] \xrightarrow{\pi} \FF[y_1,\ldots,y_m] / I,
\]
where $\phi$ is defined by $x_i \mapsto h_i(y_1,\ldots,y_m)$ and $\pi$ is the canonical map.
$\phi$ maps onto $T$.
If we knew $K = \ker \pi \circ \phi$, the image $\phi(K)$ would be precisely $T \cap I$.
Let $R = \FF[x_1,\ldots,x_n,y_1,\ldots,y_m]$.
Note that $R$ contains $T'$ and $\FF[y_1,\ldots,y_m]$ as subrings.
Thus $\phi$ can be extended to $\tilde \phi$ as
\[
 \tilde \phi: R = \FF[x_1,\ldots,x_n,y_1,\ldots,y_m] \to \FF[y_1,\ldots,y_m] \quad \text{ defined by } \quad 
\begin{cases} x_i \mapsto h_i(y), \\ y_j \mapsto y_j . \end{cases}
\]
The kernel of $\pi \circ \tilde \phi$ is the ideal $J = (x_1 - h_1(y), \ldots, x_n - h_n(y) ) + I$.
We have $K = J \cap T'$, which can be computed by the method above.
%
For example, let $I = (1+x+y+z,~1+xy+yz+zx) \subseteq \FF_2[x,y,z]$. 
If a subring is $T = \FF_2[x^3,y^3,z^3]$, then an auxiliary ring is $T' = \FF_2[x',y',z']$.
The intersection $I \cap T$ can be found by computing
\[
 (1+x+y+z,~1+xy+yz+zx,~x'-x^3,~y'-y^3,~z'-z^3) \cap \FF[x',y',z'].
\]
In the end, one has to replace $x'$ with $x^3$, $y'$ with $y^3$, and $z'$ with $z^3$.


The Gr\"obner basis is also useful to compute the vector space dimension of $R/I$, where $R=\FF[x_1,\ldots,x_n]$.
The remainder $r$ after division by a Gr\"obner basis of $I$ is unique
and has a property that $\lt(r)$ is not divisible by any leading term of the Gr\"obner basis elements.
The remainders are unique representatives for elements of $R/I$.
It follows that there is a one-to-one correspondence as sets between $R/I$ and $R/\lt(I)$.
In particular, they are isomorphic as vector spaces.
The generators of $\lt(I)$ are monomials $m_1,\ldots,m_t$ computed by the Gr\"obner basis.
Recall that monomials are in one-to-one correspondence with the hyper-octant $H$ of nonnegative coordinates of $\ZZ^n$.
Therefore, the vector space basis of $R/\lt(I)$ is labeled by points of $H$
that are not contained in any cone whose vertex, the least element with respect to the monomial order, is $m_i$ ($1\le i \le t$).
For example, let $I = (x^2,xy+y^2) \subset \FF[x,y]=R$.
The reduced Gr\"obner basis under a lexicographic order in which $x \succ y$ is $\{ y^3, xy+y^2,x^2 \}$.
Therefore, $R/\lt(I)$ has a vector space basis $\{1,x,y,y^2\}$.
In a different lexicographic order in which $x \prec y$, the reduced Gr\"obner basis of $I$ is
$\{x^2,y^2+xy\}$. A vector space basis of $R/\lt(I)$ is $\{1, x, y, xy \}$.


\subsection*{Syzygies and generalization to free modules}

So far we have discussed only ideals and generators.
It is straightforward to generalize the notions to submodules 
of free modules over polynomial rings $R = \FF[x_1,\ldots,x_n]$.
We first need to generalize the monomial order.
Let $e_1,\ldots,e_r$ be a basis of free module $F = R^r$.
A {\bf monomial} of $F$ is a product of variables times one of the basis,
e.g., $x_1 x_2^5 e_3$ and $x_1^3 x_3 e_1$ are monomials.
A {\bf monomial order} on $F$ is defined in the same way as \eqref{cond-TermOrder}.
The well-ordering property holds with the same proof.
The leading terms are defined the same way, as well as the ideal generated by leading terms.
The division algorithm does not need to be restated.
The Buchberger criterion and algorithm make sense with a minor change that
{\em if two leading terms involve different basis elements, then the S-polynomial is defined to be zero.}
The elimination order are still applicable, and the calculation techniques are valid with respect to submodules.

Nontrivial relations between polynomials are called syzygies.
(Here, polynomials include basis elements, so they are really elements of a free module.)
More formally, if a submodule is defined by the image of a module map $\phi : F' \to F$
between finitely generated {\em free} modules, the kernel of $\phi$ is called {\bf syzygy}. 
For instance, the syzygies of $x$ and $y$ can be described by a submodule $M$ of $\FF[x,y]^2$ 
generated by $\begin{pmatrix} y \\ -x \end{pmatrix}$,
since the ideal $(x,y)$ with generators $x$ and $y$ is the image of the map
$\phi : R^2 \xrightarrow{\begin{pmatrix} x & y \end{pmatrix}} R$ and $M = \ker \phi$.
The relation between $x$ and $y$ is thus $y(x) -x(y) = 0$.
We have implicitly seen how to obtain nontrivial relations from the Buchberger algorithm.
We formed an S-polynomial of two polynomials and ran the division algorithm with respect to a set $G=\{g_i\}$ of polynomials:
\[
 \sigma(g_i,g_j) = \frac{\lt(g_j)}{\gcd(\lt(g_i),\lt(g_j))} g_i - \frac{\lt(g_i)}{\gcd(\lt(g_i),\lt(g_j))} g_j = \sum_k h^{ij}_k g_k + r ,
\]
where $h^{ij}_k \in R$.
If $G$ were a Gr\"obner basis, then by Buchberger criterion we would have $r=0$, 
and the expression would be a nontrivial relation among $g_i$.
More formally, let the module be generated by a Gr\"obner basis $\{g_i\}$.
That is, module is the image of the map $\phi : F' \to F$
where the basis $\{\epsilon_i\}$ of $F'$ is mapped as $\epsilon_i \mapsto g_i$.
Then,
\[
\tau_{ij} = \frac{\lt(g_j)}{\gcd(\lt(g_i),\lt(g_j))} \epsilon_i - \frac{\lt(g_i)}{\gcd(\lt(g_i),\lt(g_j))} \epsilon_j - \sum_k h^{ij}_k \epsilon_k 
\]
is mapped to zero under $\phi$.
It is a theorem that {\em $\ker \phi$ is generated by $\tau_{ij}$.}
$\{\tau_{ij}\}$ is often a redundant generating set for $\ker \phi$.
Be warned that $\tau_{ij}$ is defined to be identically zero if $\lt(g_i)$ and $\lt(g_j)$ involve different basis elements of $F$.

Given a finite presentation of a module $F_0/\im \phi_1$ where $\phi_1 : F_1 \to F_0$ is a module map 
between finitely generated free modules,
we can algorithmically find $\ker \phi_1$. This kernel is a finitely generated submodule of $F_1$, 
so it can be identified with the image of a map $\phi_2 : F_2 \to F_1$ between finitely generated free modules.
We can continue this as many times as we want, and obtain a sequence of maps
\[
\cdots \to F_n \xrightarrow{\phi_n} F_{n-1} \to \cdots \to F_2 \xrightarrow{\phi_2} F_1 \xrightarrow{\phi_1} F_0 .
\]
This sequence is called a {\bf free resolution} of $F_0/\im \phi_1$.
A particularly interesting case is when the resolution is finite, i.e., $F_n =0$ for some $n$.



\section{Localization}
\label{sec:localization}

\subsection*{Tensor product}
Operations of modules include direct sum and quotient,
both of which are automatic from the fact that modules are abelian groups.
Here we introduce another operation on modules called tensor product.
It is perhaps best defined by a categorical characterization,
which is the most useful to prove things.
Here we define it in a colloquial way.

Let us first review a familiar case of vector spaces over a field $\FF$.
The tensor product $V \otimes W$ of two vector spaces $V,W$ is
a collection of formal finite linear combinations $\sum_{i,j} a_{ij} v_i \otimes w_j$ 
of ``products'' $v_i \otimes w_i$ where $v_i \in V,~ w_j \in W$ and $a_{ij} \in \FF$.
The ``product'' $\otimes$ is multilinear since
\[
 (v_1 + v_2) \otimes w = v_1 \otimes w + v_2 \otimes w \quad \text{and} \quad
 v \otimes (w_1 + w_2) = v \otimes w_1 + v \otimes w_2.
\]
The scalar multiplication floats around $\otimes$ as
\[
 (a v) \otimes w = v \otimes (a w)
\]
for any $a \in \FF$.
In physics literature, it is sometimes said that a tensor is a multi-indexed object
that transforms in a certain way under transformations for ``indices.''
This phrasing focuses on the coefficients $a_{ij}$ and captures the multilinearity, as
\[
 \sum_{ij} a_{ij} v_i \otimes w_j = \sum_{ij} a_{ij} \left( \sum_a A_{ia}v_a' \right) \otimes \left( \sum_b B_{jb} w_b' \right)
 = \sum_{ab} \left(\sum_{ij}a_{ij} A_{ia}B_{jb} \right) v_a' \otimes w_b'.
\]


The {\bf tensor product of two modules} is defined in a similar way.
For an $R$-module $M$ and $N$,
we build a formal abelian group out of all possible expressions $x \otimes_R y$ for $x \in M$ and $y \in N$,
with the following equalities {\em imposed}:
\begin{align}
 (x + x') \otimes_R y &= x \otimes_R y + x' \otimes_R y , \nonumber \\
 x \otimes_R (y + y')  &= x \otimes_R y + x \otimes_R y' ,\nonumber \\
 r x \otimes_R y &= x \otimes_R r y .\label{R-linear-tensor}
\end{align}
The resulting abelian group is denoted by $M \otimes_R N$.
It is an $R$-module by defining an action $r \cdot (x \otimes_R y) = rx \otimes_R y$.
This is the tensor product of $M$ and $N$.
When there is no confusion about the base ring $R$, we write $\otimes$ instead of $\otimes_R$.

Note that any abelian group $M$ is a module over $\ZZ$ by action $n \cdot m = \sum_{i=1}^{|n|} \mathrm{sgn}(n)m$
for any $n \in \ZZ$ and any $m \in M$.
Hence, we can take the tensor product of any abelian groups as $\ZZ$-modules.
Since any $R$-module is an abelian group,
one realizes that there are at least two ways to form tensor products of two $R$-modules $M$ and $N$;
$M \otimes_\ZZ N$ and $M \otimes_R N$.
All but one equations of \eqref{R-linear-tensor} remain unchanged.
The last equation declaring $R$-linearity of $\otimes_R$ depends on $R$.
It makes a huge difference.
For instance, $\mathbb{C} \otimes_\mathbb{C} \mathbb{C} \cong \mathbb{C}$ is one dimensional,
but $\mathbb{C} \otimes_\QQ \mathbb{C}$ is infinite dimensional
since there are infinitely many irrational numbers.
$\mathbb{C} \otimes_\ZZ \mathbb{C}$ is even larger.
When the subscript is omitted, {\em the tensor product of two $R$-modules is taken over $R$ by convention}.

Since the tensor product of two modules is a module,
it makes sense to take the tensor product of three or more modules.
Fortunately, the order of the tensor product does not matter.
\[
 (L \otimes_R M) \otimes_R N \cong L \otimes_R (M \otimes_R N) \quad  \quad \quad M \otimes_R N \cong N \otimes_R M
\]


There are (unwelcome) phenomena for general modules that never occur for vector spaces.
The tensor product of two nonzero modules may be zero.
For example, the tensor product of two $\ZZ$-modules $\ZZ/(2)$ and $\QQ$ is zero
because $[a] \otimes b = [a] \otimes 2\frac{a}{2} = 2[a] \otimes \frac{a}{2} = 0 \otimes \frac{a}{2} = 0$.
An expression $x \otimes y$ may be zero in $M \otimes N$ but may not be zero in $M' \otimes N'$
where $M' \le M$ and $N' \le N$ are submodules.
For example, let $M=\ZZ$ and $N = \ZZ/(2)$ be $\ZZ$-modules. Choose $M' = 2\ZZ$ and $N' = N$.
Now $2 \otimes [1]$ is nonzero in $M' \otimes N'$, 
but in $M \otimes N$, it is equal to $2 \cdot 1 \otimes [1] = 1 \otimes [2] = 0$.


Below we explain a ``safe'' and powerful tensor product.

\subsection*{Ring of fractions}

The ring of rational numbers $\QQ$ is constructed from $\ZZ$ by inverting nonzero elements.
We wish to do a similar thing for general rings.
Let $U$ be a {\bf multiplicatively closed} subset of a ring $R$,
i.e., $1 \in U$ and any product of two elements of $U$ lies in $U$.
$U$ needs not be closed under addition.
For example, the set of all nonzero numbers in $\ZZ$ is a multiplicatively closed set.
More important is the complement of a prime ideal $\pp$.
$1$ is included in $R \setminus \pp$ because $\pp \neq (1)$.
If $a \notin \pp$ and $b \notin \pp$, then $ab \notin \pp$.
This is the defining property of the prime ideal.

We construct ``fractions'' by putting elements of $U \subseteq R$ in denominators
and elements of $R$ in numerators. The set of all fractions becomes a ring with the ``usual'' additions and multiplications
\begin{align}
 \frac{a}{p} + \frac{b}{q}    &= \frac{aq + bp}{pq}, \label{fraction-add} \\
\frac{a}{p} \cdot \frac{b}{q} &= \frac{ab}{pq}. \label{fraction-multiply}
\end{align}
Moreover, the elements of $U$ are invertible!
\[
 \frac{1}{p} \cdot p = \frac{p}{p} = 1, \quad \quad \frac{q a}{qp} = \frac{a}{p}.
\]
The new {\bf ring of fractions}%
\footnote{
Rigorously, the ring of fractions is a collection of equivalence classes of $R \times U$;
$(a,p) = (b,q)$ if and only if there exists $s \in U$ such that $s(aq-bp) = 0$.
The reason we do not define the equivalence using ``$aq=bp$'' is
that there could be zero-divisors in $R$.
Interested readers might want to check that 
$U^{-1}R$ indeed becomes a ring using \eqref{fraction-add} and \eqref{fraction-multiply} as definitions.
}
is denoted by $U^{-1}R$ or $R[U^{-1}]$.
Two special cases are so important that they deserve separate notations.
\begin{itemize}
 \item If $U = R \setminus \pp$, the ring of fractions is called the {\bf localization of $R$ at $\pp$} and denoted by $R_\pp$.
 \item If $U = \{1, f, f^2, f^3, \ldots\}$ for some $f \in R$, the ring of fractions is denoted by $R_f$.
\end{itemize}
The original ring $R$ lives in $U^{-1}R$ via a canonical map $r \mapsto \frac{r}{1}$.
With this canonical map, we omit $1$ in the denominators.
Note that $U^{-1}R$ is an $R$-module.
For example, $\ZZ_{2} = \ZZ[\half]$ is a ring of fractions with denominators are powers of $2$.
(Yes, this is a confusing notation since $\ZZ_n$ is used to mean $\ZZ/(n)$.)
$\ZZ_{(2)}$ is a ring of fractions with odd denominators.
$\QQ[x]_{(x)}$ is a ring of fractions of polynomials where denominators does not vanish at $x=0$.
What is $\QQ[x,y]_{xy}$? It is a ring of fractions of polynomials where denominators are powers of $xy$.
Since $\frac{1}{x} = \frac{y}{xy}$ and $\frac{1}{y} = \frac{x}{xy}$,
we notice that $\QQ[x,y]_{xy}$ is really the ring of Laurent polynomials.


Modules can be fractionalized as well.
Let $M$ be an $R$-module.
Choose a multiplicatively closed set $U$ of the base ring $R$,
and define $U^{-1}M$ as the set of all fractions with elements of $M$ in the numerators and elements of $U$ in the denominators.
The addition within $U^{-1}M$ is defined by \eqref{fraction-add}.
$U^{-1}M$ becomes an $U^{-1}R$-module using \eqref{fraction-multiply}.
In fact, {\em $ U^{-1}M \cong (U^{-1}R) \otimes_R M $ as $U^{-1}R$-modules.} Moreover,
\begin{itemize}
\item {\em $U^{-1}M \otimes_{U^{-1}R} U^{-1}N = U^{-1}( M \otimes_R N)$ for any $R$-modules $M,N$.}
\end{itemize}




The {\bf localization}, the process passing to fractional rings and modules,
is so well behaving under nearly all conceivable operations.
Let $A, B \le M$ be submodules.
\begin{itemize}
\item {\em $U^{-1}A$ is a submodule of $U^{-1}M$}, i.e., $U^{-1}A$ injects into $U^{-1}M$.
\item {\em $U^{-1}A \cap U^{-1}B = U^{-1}(A \cap B)$}.
\item {\em $U^{-1}A + U^{-1}B = U^{-1}(A + B)$}.
\item {\em $U^{-1}(M/A) = (U^{-1}M)/(U^{-1}A)$}.
\end{itemize}
They all originate from the following.
A sequence of maps among modules $A,B,C$
\[
 0 \to A \to B \to C \to 0
\]
is a {\bf short exact sequence} if
\begin{itemize}
 \item $\ker( A \to B ) = 0$,
 \item $\im( A \to B ) = \ker ( B \to C)$,
 \item $\im(B \to C ) = C$.
\end{itemize}
{\em The localization preserves short exact sequence}, i.e., if $0 \to A \to B \to C \to 0$ is a short exact sequence, then
\[
 0 \to U^{-1}A \to U^{-1}B \to U^{-1}C \to 0
\]
is also a short exact sequence. For example, the first and the fourth statement above follows from localizing the short exact sequence
\[
 0 \to A \to M \to M/A \to 0.
\]




The localizations at prime ideals are useful because {\em an $R$-module $M$ is zero if and only if
$M_\pp = 0$ for every prime ideal $\pp$ of $R$.}
In fact, $M = 0$ {\em if and only if $M_\mm = 0$ for every maximal ideal $\mm$ of $R$.}
This can be understood as follows.
The {\bf annihilator} denoted by $\ann_R M$ of an $R$-module $M$ is an ideal of $R$ defined by
\begin{equation}
 \ann_R M = \{ r \in R ~|~ r m = 0 \text{ for any } m \in M \}
\end{equation}
It is an ideal because $(a+b)m = am + bm = 0 + 0 = 0$ and $(ra)m = r(am) = 0$ if $a, b \in \ann_R M$.
Note that {\em any $R$-module $M$ becomes a module over $R/(\ann_R M)$. }
What is to be proved for this statement is the well-definedness of the defining operation \eqref{cond-module},
i.e., one needs to check that $[a]\cdot m = [b]\cdot m$ if $[a]=[b] \in R/(\ann_R M)$ for any $m \in M$.
Note that if $R$ is a polynomial ring over a field, there is an algorithm to compute $\ann_R M$
for a finitely presented module $M$ using pull-backs and Gr\"obner basis techniques.
Now, it is easy to see that 
{\em a module $M_\pp$ (the localization of $M$ at a prime $\pp$)
is zero if and only if the annihilator is not contained in $\pp$},
since a module is zero if and only if $1$ or any unit is an annihilator,
but the localized ring at $\pp$ anything outside $\pp$ is a unit.
If $M$ becomes zero at the localization at any maximal ideal,
it means that the annihilator is not contained in any maximal ideal.
The only elements of $R$ that lie outside of any maximal ideal are units.
Therefore, $M = 0$.


Recall that there exists a canonical map $\phi$ from $R$ to a ring of fractions $U^{-1}R$, sending $r$ to $\frac{r}{1}$.
The image under $\phi$ of any ideal $I$ of $R$ is not in general an ideal.
However, we may consider the ideal of $U^{-1}R$ generated by $\phi(I)$.
This ideal is denoted by $I^e \subseteq U^{-1}R$, called an {\bf extended ideal}.
In fact, {\em any ideal of $U^{-1}R$ is an extended ideal.}
To see this, consider any ideal $J$ of $U^{-1}R$, and its {\bf contraction} $\phi^{-1}(J) \subseteq R$.
An element of $J$ is $\frac{x}{s}$. Since $J$ is an ideal, $\frac{s}{1} \frac{x}{s} = \frac{x}{1} \in J$.
That is, $x \in \phi^{-1}(J)$. We have shown $(\phi^{-1}(J))^e = J$.
Note that {\em every prime ideal of $U^{-1}R$ is an extended ideal of a prime ideal of $R$ 
that do not intersect $U$.}
We have seen that the inverse image of a prime ideal $\pp$ is prime;
the contraction $\pp^c = \phi^{-1}(\pp)$ of a prime ideal of $U^{-1}R$ is prime.
Since prime ideal is not $(1)$, it cannot contain any unit.
Hence $\pp^c$ cannot meet the set of units $U$.
Conversely, for any prime ideal $\mathfrak q$ of $R$ that do not intersect $U$,
$\mathfrak q^e$ is a prime ideal of $U^{-1}S$,
since $\frac{ab}{ss'} \in \mathfrak q^e \Longleftrightarrow \frac{ab}{1} \in \mathfrak q^e$.

It is clear that if $\phi^{-1}(J)$ is generated by $n$ elements of $R$,
then $J$ is generated by $n$ elements of $U^{-1}R$.
More importantly, {\em if $R$ is a Noetherian ring, then any localization is a Noetherian ring.}


A {\bf local ring} is a ring with a unique maximal ideal.
Typically, maximal ideals are not unique.
In the ring of integers, every prime number generates a distinct maximal ideal.
Let $R_\pp$ be a localized ring at a prime ideal $\pp$ of $R$.
By definition of $R_\pp$, the set of all denominators do not meet $\pp$.
Hence, $\pp_\pp$ is a prime ideal of $R_\pp$.
If $\mm$ is any ideal of $R_\pp$, we must have $\mm \subseteq \pp_\pp$.
Otherwise, $\mm$ contains an element outside $\pp_\pp$, which is invertible,
and hence $\mm = (1)$.
The localized ring $R_\pp$ at a prime ideal is a local ring with a unique maximal ideal $\pp_\pp$.
One thing to remember about a local ring is that
{\em any element outside the unique maximal ideal is a unit, i.e., it is invertible.}



\subsection*{Geometry}

Why is it called ``localization?''
Here we briefly introduce the algebro-geometric point of view on rings.

The (only) algebraic structure that we should start with for a geometric object is a function space.
By a function we mean a set-theoretical map from the geometric object to numbers.
We can impose certain conditions such as continuity or differentiability
on the function space to study deeper and interesting aspects of the geometric object.
In classical algebraic geometry, the geometric object, called variety,
is defined by one or more polynomial equations,
and the functions are defined by polynomials.

Roughly speaking, there are two interesting types of functions:
globally defined functions and locally defined functions.
The locally defined ones may not be well-defined for some regions far from the point of interest.
In the algebraic setting, the global functions are given by
polynomial expressions without denominators.
There may be several polynomial expressions for the same function.
This happens when two polynomials differ by the equations of the variety.
For example, the global function space, called {\bf coordinate ring}
of a parabola on $\RR^2$ defined by $y = x^2$ is a quotient ring $R = \RR[x,y]/(y-x^2)$.

Locally defined functions at a point $p$, if they are given as a fraction of two polynomials,
are precisely those whose the denominators are not zero at $p$.
For example, the fraction $\frac{1}{x+1}$ is a locally defined function at a point $(0,0)$ of the parabola,
but $\frac{x-1}{y}$ is not.
Collecting all locally defined functions at $\mm = (0,0)$, we obtain
a {\bf local ring} $R_\mm = (\RR[x,y]/(y-x^2))_{(x,y)}$,
which is consisted of fractions whose denominators are not contained in the maximal ideal $(x,y)$.
The global function ring $R$ is \emph{localized at $\mm$ to be a local ring $R_\mm$}.
Note that we have identified a point on the parabola with a maximal ideal $\mm$.
This makes sense because the set of all solutions 
of the equations $f(x,y)=0$ where $f(x,y) \in \mm$ is exactly $(0,0)$.

One might ask what it means geometrically to localize at a prime ideal $\pp$.
A prime ideal defines a subvariety by equations $f = 0$ where $f \in \pp$.
The localization reveals the set of ``functions'' that are locally defined ``around'' the subvariety.
(The notion of neighborhood can be made rigorous by defining a topology using polynomials.)




\section{Determinantal ideal}
\label{sec:det-ideal}

When is a square matrix with integer entries invertible?
A naive answer is that it is when the determinant is nonzero.
It is true when the matrix is viewed as a matrix over $\QQ$.
If we insist that the inverse matrix must be expressed over $\ZZ$,
then the answer is that it is when the determinant is $\pm 1$.
That is, the determinant must be invertible within the ring.
Put differently, it is when the ideal generated by the determinant is the unit ideal.
The determinant has another use. The rank of a matrix is the largest $k$ such that
there exists a $k \times k$ submatrix whose determinant is a nonzero.
Put differently, it is when the ideal generated by determinants of all $k \times k$ submatrices
is nonzero.

These observations motivate us to define a determinantal ideal over an arbitrary commutative ring $R$.
Consider a rectangular matrix $\mathbf M$ with entries in $R$.
A {\bf $k^\text{th}$ minor} is the determinant of a $k \times k$ submatrix.
(There are quite many $k^\text{th}$ minors if $k$ is about half of the matrix size.)
Note that to define the minor we do not have to assume that the full matrix $\mathbf M$ is square.
When $k$ is larger than the number of rows or columns, $k^\text{th}$ minor is, by definition, zero.
{\bf $k^\text{th}$ determinantal ideal $I_k(\mathbf M)$} of $\mathbf M$ is the ideal of $R$ 
generated by all $k^\text{th}$ minors of $\mathbf M$.
By convention, $0^\text{th}$ determinantal ideal is taken as $I_0(\mathbf M) = (1) = R$.
If $\mathbf M$ is $n \times n$, then $n^\text{th}$ determinantal ideal is generated by a single element,
the determinant of $\mathbf M$.
The first determinantal ideal is generated by all entries of $\mathbf M$.
A $(k+1)^\text{st}$ minor is a linear combination of $k^\text{th}$ minors,
as it has a cofactor expansion.
Therefore, the determinantal ideals form a decreasing chain:
\[
 R = I_0(\mathbf M) \supseteq I_1(\mathbf M) \supseteq I_2(\mathbf M) \supseteq \cdots 
    \supseteq I_{\min( \# row, \# col)}( \mathbf M ) \supseteq (0)
\]
The {\bf rank} of $\mathbf M$ is the largest $k$ such that $I_k(\mathbf M)$ is nonzero.


The determinantal ideals are invariants of $\mathbf M$ 
under invertible matrix multiplications on the right or on the left.
\[
 I_k ( \mathbf {M D} ) = I_k( \mathbf{E M} ) = I_k (\mathbf M) \quad 
\quad \text{if $\mathbf{D,~E}$ are square and invertible.}
\]
To understand this, it is enough to see $I_k( \mathbf{E M} ) \subseteq I_k( \mathbf M )$ for arbitrary matrix $E$.
Then, for invertible $\mathbf E$, we would have $I_k(\mathbf{E^{-1} E M }) \subseteq I_k(\mathbf{E M})$.
We need to see that a $k^\text{th}$ minor of $ \mathbf{E M}$ is a linear combination of $k^\text{th}$ minors of $\mathbf M$.
Fix a $k \times k$ submatrix $\mathbf K$ of $\mathbf{EM}$.
\[
 \mathbf K_{il} = \sum_{j} e _{ij} \mathbf M _{j l} 
\]
where $e_{ij} = \mathbf{E}_{ij}$ is written with small letter to emphasize it is merely an element of $R$.
Now the determinant of $\mathbf K$ is vividly expressed by minors of $\mathbf M$,
since the determinant is multilinear over $R$ in rows $K_i$ of $\mathbf K$.
\begin{align*}
 \det \mathbf K 
&= \det( K_1, \ldots, K_k ) 
= \det \left[ \sum_{j_1} e_{1 j_1} M_{j_1}, \ldots, \sum_{j_k} e_{k j_k} M_{j_k} \right]\\
&= \sum_{j_1,\ldots,j_k} e_{1 j_1} \cdots e_{k j_k} ~\det \left[  M_{j_1},\ldots, M_{j_k} \right]
\end{align*}
where $M_j$ is the row $j$ of the submatrix of $\mathbf M$ that is consisted of entries contributing to $\mathbf K$.


Suppose $R$ be a local ring with a maximal ideal $\mm$,
and consider a matrix $\mathbf M$ over $R$ of rank $k$.
If $I_k(\mathbf M) = (1)$, how would $\mathbf M$ look like after row and column operations?
The operations are invertible, and therefore do not change the determinantal ideals.
Since the determinantal ideals form a descending chain,
it follows that $I_1(\mathbf M) = (1)$.
It means that at least one entry must lie outside $\mm$ and is invertible.
Bringing that entry to $(1,1)$ position by permuting rows and columns,
we can eliminate other entries in the first row and column, by row and column additions.
So
\[
 \mathbf M \cong \begin{pmatrix} 1 & 0 \\ 0 & \mathbf M' \end{pmatrix}
\]
where $\cong$ means the equivalence up to invertible matrix multiplication on the left or right.
Now we use the fact that $I_2(\mathbf M) = (1) = I_1( \mathbf M' )$.
By a similar process we can extract $1$ from $\mathbf M'$.
After $k$ steps, we have
\begin{equation}
 \mathbf M \cong \begin{pmatrix} \id_k & 0 \\ 0 & 0 \end{pmatrix}.
\label{free-module-presentation}
\end{equation}
Here, observe that $R$ is not necessarily a field.
It is just a local ring where some nonzero elements may not be invertible.
The smallest nonzero determinantal ideal being a unit ideal
was strong enough to imply the structure of the matrix.

\subsection*{Smith normal form}

Beyond the ``easy'' local ring case,
there is one more easy case of {\bf principal ideal domains},
in which every ideal is generated by a single element.
It is not necessarily local.
{\em The ring of integers and the ring of polynomials in one variable over a field
are principal ideal domains};
given finitely many generators of an ideal,
one can find the greatest common divisor ($\gcd$) by the Euclidean algorithm.

Let $R$ be a principal ideal domain.
We ask what normal forms of matrices we may have, after
invertible matrix multiplication on the left or right.
In other words, we seek for invariants of the matrices.
Observe that the $\gcd$ of two elements%
%%%%
\footnote{To speak of $\gcd$, it must be first proved that in $R$ any element has a unique factorization.
It is true that in any principal ideal domain, any element is a product of irreducible factors
that are unique up to units. If the latter is satisfied, the ring is called a {\bf unique factorization domain}.
Note that an arbitrary ring may not be a unique factorization domain.
For example, in $\ZZ[\sqrt{-5}]$, we have $6 = 2 \cdot 3 = (1+\sqrt{-5})(1-\sqrt{-5})$,
which are two different factorizations of $6$ into irreducible factors.}
%%%%
$f$ and $g$ can be expressed as a linear combination of $f$ and $g$
\[
\gcd(f,g) = a f + b g
\]
for some $a,b \in R$. To see this, consider the ideal $(f,g)$.
It must be generated by a single element, say, $d$, i.e., $(d) = (f,g)$.
Then, $f \in (d) \Longleftrightarrow f = d f'$ and $g \in (d) \Longleftrightarrow g = d g'$.
So $d$ is a common divisor of $f$ and $g$. In addition, $(d) = (f,g) \subseteq (\gcd(f,g))$
implies that $\gcd(f,g)$ divides $d$. Thus, $\gcd(f,g)$ and $d$ are the same up to units,
and $\gcd(f,g) \in (f,g)$. We thus obtain the above expression.
Now, consider the following matrix equation
\[
\begin{pmatrix} \gcd(f,g) \\ 0 \end{pmatrix}
= \begin{pmatrix} a & b \\ -g' & f' \end{pmatrix} \begin{pmatrix} f \\ g \end{pmatrix}
\]
The $2 \times 2$ matrix has determinant $af' + bg' = \gcd(f,g)/d$, a unit.
We have transformed a $2 \times 1$ matrix by an invertible left multiplication
such that there is only a single nonzero entry left. 
A similar thing can be done for an $n \times 1$ matrix.
\[
\begin{pmatrix} \gcd(f_1,\ldots,f_n) & 0 & \cdots & 0 \end{pmatrix}^T
= \mathbf E \begin{pmatrix} f_1 & \cdots & f_n \end{pmatrix}^T
\]
where $\mathbf E$ is invertible.

If we are given with a rectangular matrix $\mathbf M$,
we can perform this transformation by looking at the first column of $\mathbf M$.
The transformed matrix $M'$ will have a nonzero entry on the first row in the first column.
Perform a similar transformation by right multiplication, focusing on the first row of $\mathbf M'$.
The new matrix $\mathbf M''$ has unique nonzero entry at $\mathbf M''_{1,1}$ in the first row,
but the first column may be screwed up.
We can iterate these transformations as many times as we want.
Will this process eventually terminate? Yes.
For example, if we are working in the ring of integers,
the absolute value of the entry $\mathbf M''_{1,1}$ is smaller than $|\mathbf M'_{1,1}|$, which is $\le |\mathbf M_{1,1}|$.
Positive integers cannot decrease forever, and the process must terminate.
More generally, the ascending chain condition makes the proof smooth.
The ideals generated by the $(1,1)$-entries will form an increasing sequence of ideals.
Our ring $R$, being a principal ideal domain, is Noetherian.
The sequence must become stationary after finitely many iterations.
If one is not too familiar with Noetherian rings,
one can rely on the fact that there are finitely many factors of a given element.
If $(d_1) \subseteq (d_2) \subseteq \cdots$ is an increasing sequence of ideals,
we have $d_{i+1} | d_{i}$. Since there are only finitely many factors in $d_1$,
one cannot have an infinite and strictly increasing sequence of ideals.

The matrix $\mathbf M^{(n)}$ obtained from $\mathbf M$ by iterating
the $\gcd$-computations $n$-times, has a unique nonzero entry $d$ at $(1,1)$
in its row and column.
\[
\mathbf M^{(n)} = \begin{pmatrix} d & 0 \\ 0 & M_1 \end{pmatrix}
\]
We can claim a bit more. If there is any entry at $(u,v)$ of $M_1$ that is not divisible by $d$,
we can add the row $u$ to the first row and run the $\gcd$-computation again.
The ideal generated by $(1,1)$-entry will become larger.
Iterating, we eventually find the largest possible ideal generated by $(1,1)$-entry.
Therefore, we can obtain $M_1$ whose entries are all divisible by $d$.

The above algorithm applied to $M_1$, and to its submatrix, and so on,
will produce a diagonal matrix. This diagonal matrix is called the {\bf Smith normal form} of $\mathbf M$.
The nonzero diagonal entries $d_1, \ldots, d_k$ have a property that
\[
d_1 | d_2 | \cdots | d_k .
\]
Here, the number $k$ is precisely the rank of the matrix $\mathbf M$.
Moreover, the determinantal ideals are 
\[
I_t ( \mathbf M ) = (d_1 \cdots d_t).
\]
This immediately proves that $d_i$'s are uniquely determined by $\mathbf M$.
The invariants $d_i$'s are called {\bf elementary divisors} of $\mathbf M$.
Since any matrix, not necessarily square, can be brought to the Smith normal form
by invertible transformations,
it follows that the elementary divisors are {\bf complete invariants},
i.e., the elementary divisors are the same for two matrices if and only if two matrices are related by
invertible matrix multiplication on the left and right.


\subsection*{Finitely generated modules and Fitting ideals}

We briefly noted about finite presentations of modules when we discussed free modules.
We say that an $R$-module $M$ is finitely presented by a matrix $\phi : R^m \to R^n$
if $M$ is isomorphic to $R^n / \im \phi$. This is more commonly denoted as $ M = \coker \phi$.

To get some feeling, let us consider a simple case where $R = \ZZ$.
The matrix $\phi: R^m \to R^n$ is an $n \times m$ matrix with integer entries.
Any basis change in $R^n$ amounts to an invertible left multiplication on $\phi$.
Any basis change in $R^m$ amounts to an invertible right multiplication on $\phi$.
We know a very convenient canonical form of $\phi$ --- the Smith normal form.
Let $\phi$ be in the Smith normal form with the diagonal elements $d_1, \ldots, d_k$.
In the simplest case $k = m = n =1$, the module $M$ is $R / \im \phi = \ZZ / (d_1)$.
If $m = n = 2$ and $d_2 = 0$, the module $M$ is $R^2 / \im \phi = \ZZ/(d_1) \oplus \ZZ$.

In fact, we can prove the structure theorem for finitely generated abelian groups very easily.
Let $M$ be a finitely generated abelian group. It can be viewed as a finitely generated
$\ZZ$-module. Let $n$ be the number of generators. We have a module map $\tilde \phi : \ZZ^n \to M$,
which is surjective. The kernel $K$ is also finitely generated since $\ZZ^n$ is Noetherian.
Hence, we have another map $\phi : \ZZ^m \to K \subseteq \ZZ^n$ where $m$ is the number of
generators of $K$. We have constructed a finite presentation of $M$ by $\phi$, i.e., $M = \ZZ^n / \im \phi$.
Bring $\phi$ to the Smith normal form by basis changes in $\ZZ^n$ and $\ZZ^m$.
If $d_1 | \cdots | d_k$ are elementary divisors of $\phi$,
we define $d_{k+1} = \cdots = d_n = 0$.
Then we have an isomorphism of $\ZZ$-modules $M \cong \ZZ/(d_1) \oplus \cdots \oplus \ZZ/(d_n)$.
Viwed as a group, the notation will be
\[
 M = \ZZ/d_1 \ZZ \times \cdots \times \ZZ/d_k \ZZ \times \ZZ^{n-k} .
\]
This is the most general form of a finitely generated abelian group.
Note that the rank $n-k$ and the elementary divisors 
$d_1 | d_2 | \cdots | d_k \neq 0$ are uniquely determined by the group $M$.
Any finite abelian group is obviously finitely generated.


Let $R$ be an arbitrary commutative ring, and $M$ be an $R$-module 
with a finite presentation given by a matrix $A: R^m \to R^n$, i.e., $M \cong R^n/\im A = \coker A$.
We define {\bf Fitting ideals} $F_i(M)$ of $M$ as
\begin{equation}
 F_i(M) = I_{n-i}(A),
\label{Fitting-ideal}
\end{equation}
so
\[
 F_0(M) \subseteq F_1(M) \subseteq \cdots \subseteq F_n(M) = R .
\]
One should ask: $A$ is one of many possible finite presentations.
How can we be sure that $F_i(M)$ is determined by $M$ and is independent of a particular presentation $A$?
We can be sure.
The presentation really consists of two sets of things:
a set of generators $x_1,\ldots,x_n \in M$ and their relations (syzygies) given by columns of $A$
\[
 \sum_{i} x_i A_{ij} = 0.
\]
Thus, we can think of $A$ as the matrix of relations.
The isomorphism $M \cong \coker A$ ensures that any possible relations among $x_i$ is generated by columns of $A$.
That is, for any coefficients $c_i$ such that $\sum_i x_i c_i = 0$ we can express $c_i$ as $c_i = \sum_j  A_{ij} d_j$.
Therefore, if we write a $n \times \infty$ matrix $\tilde A$ by collecting \emph{all} relations among $x_i$,
then the $k^\text{th}$ determinantal ideal of $\tilde A$ is exactly $I_k(A)$.
In other words, $I_k(A)$ is determined by the chosen generators of $M$.
Our question on the well-definedness of the Fitting ideal concerns, in fact, many possible choices of generators for $M$.
%%%
Now, let $y_1, \ldots, y_{n'}$ be elements of $M$.
We have $n+n'$ generators $x_i, y_{i'}$ of $M$,
and $M$ is presented as a quotient module of $R^{n+n'}$. 
The relations among $x_i,y_{i'}$ constitute a matrix
\[
W =
\begin{pmatrix}
 A & A'  & \# \\
 0 & \id & \# \\
\end{pmatrix} 
\sim
W' =
\begin{pmatrix}
 A & 0   & A' \\
 0 & \id & 0 \\
\end{pmatrix} 
\sim
W'' =
\begin{pmatrix}
 A & 0   & 0 \\
 0 & \id & 0 \\
\end{pmatrix} 
\]
where $\sim$ means equality up to row or column operations.
The first matrix $W$ can be written as shown because $y_{i'}$ can be written as some combination of $x_i$.
$W''$ is obtained from $W'$ since we know that $A$ generates all relations among $x_i$.
Since determinantal ideals are invariant under invertible matrix multiplications,
we see that $I_{k+n'}(W) = I_{k+n'}(W') = I_{k+n'}(W'') = I_k(A)$.
Another finite presentation $B : R^{m'} \to R^{n'}$ of $M \cong \coker B$
gives another set of generators.
By the above calculation, $I_{n-k}(A) = I_{n+n'-k}(W) = I_{n'-k}(B)$.
Therefore, {\em the Fitting ideals are well-defined.}
We now see clearly why the numbering of the Fitting ideals are given as \eqref{Fitting-ideal}.
(One might notice that the finite presentation is actually too much than 
what is needed to define Fitting ideals. They can be defined for any finitely generated module.)


The first nonvanishing Fitting ideal is important because it tells 
when a finitely generated module becomes free after localization.
For example, suppose an $R$-module $M$ is finitely presented by $\phi : R^m \to R^3$ and $F_1(M)$ is nonzero.
Let $\pp$ be a prime ideal of $R$ such that $F_1(M) \not\subseteq \pp$.
Localizing at $\pp$, we see that $F_1(M) = I_2(\phi)$ becomes the unit ideal,
since anything outside $\pp$ is a unit in $R_\pp$.
Then, we have seen in \eqref{free-module-presentation} that
the matrix $\phi_\pp$ is equivalent to a diagonal matrix with entries $0$ or $1$.
Therefore, $\coker \phi_\pp$ is obviously isomorphic to $R_\pp^3/R_\pp^2 = R_\pp^1$, a free module.
Conversely, if $\coker \phi_\pp$ is free for some prime ideal $\pp$,
then $I_k(\phi_\pp)$ is either $(1)$ or $(0)$.
Therefore, every Fitting ideal of $\coker \phi_\pp$ is either $(0)$ or $(1)$.
In conclusion, {\em the localizaed module $M_\pp$ of $M$ at a prime ideal $\pp$ is free 
if and only if the first nonvanishing Fitting ideal of $M$ is not contained in $\pp$.}


The initial Fitting ideal $F_0(M)$ is also interesting because it approximates the annihilator of $M$.
Let $M$ be generated by $n$ elements. Then,
\begin{equation}
(\ann M)^n \subseteq F_0(M) \subseteq \ann M .
\label{Fitting-annihilator}
\end{equation}
Here, the {\bf ideal power} $(\ann M)^n$ means the ideal generated by all products of $n$ elements from $\ann M$.
By definition, $F_0(M) = I_n(A)$ where $A$ is a matrix of relations among the $n$ generators $x_i$ of $M$.
If $Z$ is an $n \times n$ submatrix of $A$, we have $\sum_i  x_i Z_{ij} = 0$.
Multiplying the adjugate matrix of $Z$, we see $(\det Z) x_i = 0$ for all $1 \le i \le n$.
Hence, $\det Z \in \ann M$. Since $F_0(M)$ is generated by these $\det Z$, we have the second inclusion
of \eqref{Fitting-annihilator}.
If $a_1,\ldots, a_n \in \ann M$, then the diagonal matrix made of $a_i$ expresses relations among $x_i$.
The determinant $a_1\cdots a_n$ therefore belongs to $F_0(M)$.
This proves the first inclusion of \eqref{Fitting-annihilator}.




\section{Finite fields}

A field is a nonzero commutative ring where every nonzero elements are invertible.
There is only one proper ideal, the zero ideal $(0)$.
A {\bf finite field} is a field with finitely many elements.
If there are $q$ elements in the field, we denote the field by $\FF_q$.
The minimum possible number is of course $q=2$ because we need $0 \neq 1$.
In this case it is called {\bf bianary field} $\FF_2$.
It may seem quite modest to require for a field to have finitely many elements,
but, unlike finite groups, finite fields are particularly simple ---
{\em there is a unique field of $q$ elements up to isomorphisms.}
Let us see why.


Since any ring is an additive group, $\FF_q$ is a finite abelian group.
How does it look like as an additive group?
Applying the ``structure theorem for finitely generated abelian group,''
we know $\FF_q \cong \bigoplus_{i=1}^k \ZZ/(d_i)$.
where $1 < d_1 | d_2 | \cdots | d_k \neq 0$.
It follows that for any element $x \in \FF_q$ we have $d_k x = 0$.
Here, $d_k$ is the smallest number such that $d_k x = 0$ for any $x \in \FF_q$.
If $d_k$ is not a prime number, say, $d_k = pp'$ with $p, p' < d_k$,
then $p = p\cdot 1 \neq 0$ in $\FF_q$ is invertible.
Hence, $p' = p' \cdot 1 = 0$, which is a contradiction to the minimality of $d_k$.
Therefore, $d_k$ is prime, and $d_1 = d_2 = \cdots = d_k = p$.
It follows that $| \FF_q | = q = d_1 d_2 \cdots d_k = p^k$.
This conclusion is often phrased as {\em a finite field has {\bf characteristic} $p$ for some prime number $p$
and the total number of elements is a power of $p$.}
Note that the image of $\ZZ \to \FF_q$, as there is a unique ring homomorphism from $\ZZ$ to any ring,
is $\FF_p$, a subfield of $\FF_q$.
It is of great importance that {\em $\FF_q$ is a $k$-dimensional vector space over $\FF_p$.}


Next, how does the multiplicative group of all nonzero elements of $\FF_q$ look like?
It is again a finitely generated abelian group.
Therefore, its group structure is uniquely determined by elementary divisors $1< e_1 | e_2 | \cdots |e_m \neq 0$.
In a similar logic as above, we have $x^{e_m} = 1$ for any nonzero $x \in \FF_q$.
It is an equation in the field, so it can be rewritten as $x^{e_m} -1 = 0$.
How many solutions can an equation have? At most the degree.
It means that $q-1 = e_1 e_2 \cdots e_m \le e_m$.
The only possibility is that $m = 1$ and $e_1 = q-1$.
That is, the multiplicative group $\FF_q^\times$ consisted of all nonzero elements
is isomorphic to the additive group $\ZZ/(q-1)$, a {\bf cyclic group} generated by a single element.
In other words, there exists an element $x$ in $\FF_q$ such that
$\{ x^n | n \in \ZZ \}$ is equal to the set of all nonzero elements.
Such an element $x$ is called {\bf primitive element} of $\FF_q$.
There could be many primitive elements in $\FF_q$.


Finally, let us mix the two operations, the addition and multiplication, and consider the ring structure of $\FF_q$.
Before we start analyzing $\FF_q$,  remark that for
any field $\FF$, not necessarily finite, and an irreducible polynomial $f(x)$ over $\FF$,
we can construct a larger field $\FF[x]/(f(x))$ that contains $\FF$.
A larger field is called an {\bf extension field} over the smaller field.


We noted above that $\FF_q$ is a $k$-dimensional vector space over $\FF_p$.
Fix a primitive element $\alpha$ of $\FF_q$.
Then, the set of ``vectors''  $\{1,\alpha,\alpha^2,\ldots,\alpha^k\}$ contains $k+1$ vectors,
so it cannot be linearly independent over $\FF_p$.
Consider the set $I_\alpha$ of all polynomials $f(t) \in \FF_p[t] \subseteq \FF_q[t]$ such that $f(\alpha) = 0$.
It is clearly a nonzero ideal, and therefore is generated by a single element $f_\alpha(t)$,
which divides any element in $I_\alpha$.
$f_\alpha(t)$ is unique if we demand the leading coefficient to be $1$.
A polynomial with the leading coefficient $1$ is called {\bf monic}.
The minimality of $f_\alpha(t)$ forces it to be irreducible.
$f_\alpha(t)$ is uniquely determined by $\alpha$, 
called the {\bf minimal polynomial} of $\alpha \in \FF_q$ over $\FF_p$.
We have established a ring homorphism $\FF_p[t] / (f_\alpha(t)) \to \FF_q$ such that $t \mapsto \alpha$.
Since $\alpha$ is a primitive element, this map is surjective.
Moreover, it is injective because of the choice of $f_\alpha(t)$.
It follows that the degree of $f(t)$ is actually equal to $k$,
the dimension of $\FF_q$ as an $\FF_p$-vector space,
and {\em the ring $\FF_q$ is isomorphic to a quotient ring of the polynomial ring $\FF_p[t]$}.


Recall the polynomial $h(x) = x^q -x$ becomes zero at any element of $\FF_q$.
It means that in the ring $\FF_q[x]$,
the polynomial $h(x)$ factorized into linear factors as
\begin{equation}
h(x) = x(x-\alpha)(x-\alpha^2)(x-\alpha^3) \cdots (x- \alpha^{q-1}) .
\label{h-factors}
\end{equation}
Since $x^q-x \in I_\alpha$, it follows that $f_\alpha(x) \in \FF_p[x]$ divides $x^q-x \in \FF_p[x]$.
Therefore, $f_\alpha(x) \in \FF_q[x]$ factorizes into linear factors, too.
That is, any root of $f_\alpha(x)$ can be found in $\FF_q$, not only $\alpha$.
Furthermore, {\em any root of $f_\alpha(x)$ is a primitive element.}
Let $\beta \in \FF_q$ be any root of $f_\alpha(x)$.
The ideal $I_\beta \subseteq \FF_p[x]$ contains $f_\alpha(x)$.
Hence, the minimal polynomial $f_1(x) \in I_\beta$ divides $f_\alpha(x)$.
Since $f_\alpha(x)$ is irreducible in $\FF_p[x]$, we must have $f_\alpha(x) = f_\beta(x)$.
It means that $\FF_q \ni \alpha \mapsto \beta \in \FF_q$ is an isomorphism,
guaranteeing that $\beta$ is a primitive element.


The quotient ring presentation of $\FF_q$ is not unique because it depends on $f_\alpha(t)$ which
is determined by a primitive element. 
It is thus a relevant question whether two finite fields with the same cardinality $\FF_q$ and $\FF_q'$ 
are isomorphic as rings.
Let $\alpha$ be a primitive element of $\FF_q$ with the minimal polynomial $f(x)$ over $\FF_p$,
and $\beta$ be a  primitive element of $\FF_q'$ with the minimal polynomial $g(x)$ over $\FF_p$.
We know that $f(x) | h(x)$ and $g(x) | h(x)$ where $h(x) = x^q -x$.
It follows from \eqref{h-factors} that $f(\beta^n) = 0$ in $\FF_q'$ for some $n$.
Given such $n$, we can define a ring homomorphism $\FF_p[x]/(f(x)) \ni x \mapsto y^n \in \FF_p[y]/(g(y))$.
It amounts to a ring homomorphism $\FF_q \ni \alpha \mapsto \beta^n \in \FF_q'$.
Note that {\em any nonzero ring homomorphism from a field is injective}
simply because the kernel is a proper ideal.
Since both the domain and the target are finite dimensional vector spaces,
our homomorphism is bijective, and we obtain a field-isomorphism.
In conclusion, {\em any finite field is uniquely determined up to isomorphisms by its cardinality.}


Let $f(x) \in \FF_q[x]$ be any irreducible polynomial over $\FF_q$.
We may consider an extension field $\mathbb E = \FF_q[x]/(f(x))$ over $\FF_q$.
It is still a finite field for being a finite dimensional vector space over a finite field.
In $\mathbb E$, $f(x)$ factorizes into linear factors by the same reasoning as above.
More generally, since any polynomial $g(x)$ is a product of irreducible polynomials,
by extending the field, one can factor further some of the irreducible factors.
Since there are finitely many irreducible factors, one eventually reaches an extension field
where $g(x)$ factorizes into linear factors completely, after finitely many extensions.
If we started with a finite field, then the ultimate field will still be finite of cardinality, say, $q'$.
It follows that $g(x)$ is a factor of $(x^{q'}-x)^n$ for some $n$,
where $n$ is to take care of potential multiplicity in $g(x)$.
$n$ can be chosen to be large. If one wishes, it may be of form $p^m$ 
where $p$ is the characteristic of the field.
Then, $(x^q -x)^{p^m} = x^{q' p^m} - x^{p^m}$.
Summarizing, {\em any nonzero polynomial $g(x)$ over a finite field divides $x^{p^{m'}} - x^{p^m}$
for some $m' > m$.}
